\section{Conclusion}
\label{sec:conclusion}

We reimplemented the Zhang~2016 classification CNN, fine-tuned it on COCO~2017,
and compared three deep-learning colorization paradigms on a single
leakage-controlled $1{,}000$-image \texttt{test2017} benchmark, with every
metric reported as \texttt{mean [95\% CI]}. Our fine-tuned model reaches PSNR
$23.29$\,dB, SSIM $0.919$, and LPIPS $0.192$, matching the SOTA-class DeOldify
on structure, trailing it modestly on pixel fidelity and perceptual quality, and
beating it by roughly $2.8\times$ on inference speed.

A fourth, diffusion paradigm was scoped into the original plan and then
explicitly excluded once we established that the only openly-available
colorization checkpoints depended on the deprecated Stable~Diffusion~2.1
backbone (Sec.~\ref{sec:experiments:no-diffusion}); reporting a degraded
substitute would have been misleading. The exclusion is itself a finding about
the open-source ecosystem's exposure to single-backbone deprecation
(Sec.~\ref{sec:discussion}).

Our practical recommendations follow directly:
\begin{itemize}
  \item \textbf{Automatic, latency-sensitive colorization} --- use a
        well-trained classification CNN; it is fast, stable, and close to GAN
        quality.
  \item \textbf{Maximum perceptual quality / photo restoration} --- use a GAN
        such as DeOldify, accepting the higher latency.
  \item \textbf{Human-in-the-loop control} --- use an interactive CNN
        \emph{with} hints; do not expect it to perform well unguided.
\end{itemize}

Two avenues would strengthen the study. First, re-deriving the $313$-bin gamut
hull would let us load the official ECCV head and report a faithful pretrained
baseline, isolating the true contribution of fine-tuning. Second, the diffusion
paradigm should be revisited once a maintained colorization checkpoint exists
on a still-supported backbone, completing the four-paradigm comparison with a
measured rather than inferred frontier point.
